\documentclass[11pt]{article}

\usepackage[T1]{fontenc}
\usepackage{lmodern}
\usepackage[margin=1in]{geometry}
\usepackage{booktabs}
\usepackage[numbers]{natbib}
\usepackage{url}
\usepackage[hidelinks]{hyperref}

\title{The Whole Life of a Paper, Shown on an Example}
\author{Anonymous Authors}
\date{}

\begin{document}
\maketitle

\begin{abstract}
This paper is an example that lives in your meertex account, so that you can
try everything on a paper that already has a past. It has a venue with
deadlines, a version that was handed in, a review with open threads and a
bibliography with mistakes in it. Each section explains one part of meertex
and says where to find it. Change whatever you like, because nothing here is
real. When you are done, delete the paper on its Settings tab.
\end{abstract}

\section{Introduction}\label{sec:intro}

A paper is more than its \texttt{.tex} files. It has a venue, deadlines in a
time zone nobody lives in, versions that went out, reviews from coauthors and a
bibliography that somebody has to check. meertex keeps all of it next to the
\LaTeX{} source~\cite{lamport1994latex}.

Your coauthors on this paper, Jonas Lindqvist and Amara Okafor, are made up.
They never get a mail, so write to them as much as you like.
Section~\ref{sec:writing} is about writing, Section~\ref{sec:dates} about
deadlines, Section~\ref{sec:versions} about versions,
Section~\ref{sec:reviews} about reviews and Section~\ref{sec:references} about
the bibliography.

\section{Writing Together}\label{sec:writing}

The Editor tab shows the files on the left, the source in the middle and the
PDF on the right. Everybody who has the paper open edits the same document.
You see their cursors with their names, and edits merge character by
character, because the document is a conflict-free replicated data
type~\cite{shapiro2011crdt}. meertex builds it with
Yjs~\cite{nicolaescu2016yjs,yjs}. Shared editing is much older than
that~\cite{ellis1989concurrency}, and seeing where the others work is what
keeps a shared document calm~\cite{dourish1992awareness}.

The PDF compiles two seconds after you stop typing, or at once with
Ctrl+Enter. A double click in the PDF jumps to its line in the source, and
back. The editor also underlines spelling mistakes, such as the one we left in
this sentense on purpose.

You do not have to write in the browser. The command line client keeps a
folder on your computer in step with the paper, so you can use your own editor
and your own tools~\cite{kleppmann2019local}.

\begin{verbatim}
pip install meertex
meertex login
meertex checkout example-paper
\end{verbatim}

\section{Deadlines and the Timeline}\label{sec:dates}

Most of a paper is written in the last days before its
deadline~\cite{hartmann2023deadline}, so meertex keeps the deadlines in view.
The Overview tab shows the venue, the Workshop on Example Papers, with its
dates. They are in AoE, as most calls for papers state them. The submission
deadline was extended by four days, and the overview still shows where it was
before. The next date, the notification, counts down at the top of the page.

When a date passes, meertex asks what happened. The rebuttal deadline is over,
so the overview asks whether we sent a rebuttal. Answering yes freezes the
text as it was at the deadline and calls it the version Rebuttal. The Timeline
tab shows all your papers on one chart, with the time you need to write each
of them.

\section{Versions and History}\label{sec:versions}

The text that went to the venue must never change. When we answered yes to
``Did you hand it in?'', meertex froze the files as they were at the deadline.
That is the version Submission. The version arXiv came a week later and adds
our names, which a double-blind venue must not see. The Versions tab compiles a
PDF for each of them and compares any two, as a word diff or as a PDF that
marks every change.

Between versions, meertex takes a snapshot a minute after the last edit and
keeps it in a git repository, as good research software does
too~\cite{wilson2014best}. The History tab lists the snapshots, and every
coauthor can clone the repository for a backup.

\section{Reviews}\label{sec:reviews}

After the submission, Amara reviewed it like a pull request. She commented on
passages of the text and asked for changes. Jonas reworked the text and
resolved her threads, so each thread shows the request, the answer and the
change side by side. On the Reviews tab, one thread is accepted, one waits for
its acceptance, one carries a suggestion that you can apply with one click and
one question is still open.

Jonas has also asked you to review the current text. Start your review on the
Reviews tab, select a passage and comment on it. Typing @ in a comment
completes a coauthor's name and tells them by mail.

\section{References}\label{sec:references}

We built the bibliography with mistakes in it on purpose, the kind that
reviewers find and language models make up. The References tab looks every
entry up at doi.org, Crossref, Semantic Scholar, DBLP, arXiv and Open Library,
fetches its links and lists what it finds. Table~\ref{tab:refs} says what to
expect.

\begin{table}[ht]
  \centering
  \begin{tabular}{@{}ll@{}}
    \toprule
    Entry & What is wrong with it \\
    \midrule
    \cite{hartmann2023deadline} & Made up, no index knows the paper \\
    \cite{dourish1992awareness} & A typo in the DOI \\
    \cite{nicolaescu2016yjs} & The wrong year \\
    \cite{kleppmann2019local} & The title lacks its subtitle \\
    \cite{ellis1989concurrency} & A placeholder instead of the pages \\
    \cite{collabtex} & A link to a repository that is gone \\
    \cite{wilson2014best} & Changed after somebody checked it \\
    \cite{mittelbach2004companion} & No entry at all \\
    \bottomrule
  \end{tabular}
  \caption{The mistakes in this paper's bibliography.}\label{tab:refs}
\end{table}

Where the record knows better, ``Take this value'' writes its value into the
entry. ``Add a reference'' makes an entry from a DOI, an arXiv id, an ISBN or a
link, which is how the missing entry can be made. Once you have checked an
entry by hand, ``Mark as checked'' writes a mark into the entry itself, and the
mark breaks when somebody changes the entry later.

\section{What Next}\label{sec:next}

Start your own paper with New paper on the paper list. It can start empty, from
a \texttt{main.tex} with a \texttt{refs.bib}, or from a publisher's template
zip. When this example has served its purpose, delete it on its Settings tab.

\bibliographystyle{plainnat}
\bibliography{refs}

\end{document}
